% ============================================================================
%  KCGL ERP - PMS (Performance Management System)
%  Module documentation for the client.
%  Template: same layout as the MOM Tracking document (MOM.pdf)
%  Compile with pdfLaTeX on Overleaf. Screenshots are optional - a missing
%  image prints a placeholder, so the document always compiles.
% ============================================================================
\documentclass[11pt,a4paper]{article}

\usepackage[utf8]{inputenc}
\usepackage[T1]{fontenc}
\usepackage[a4paper,margin=1in,headheight=15pt,headsep=18pt]{geometry}
\usepackage{fancyhdr}
\usepackage{graphicx}
\usepackage{float}
\usepackage{booktabs}
\usepackage{tabularx}
\usepackage{longtable}
\usepackage{array}
\usepackage{enumitem}
\usepackage{xcolor}
\usepackage{amssymb}
\usepackage{caption}
\usepackage[hidelinks]{hyperref}

% ---- page furniture --------------------------------------------------------
\definecolor{kcglblue}{HTML}{2F3192}
\definecolor{swBand}{HTML}{2F3192}
\definecolor{swHead}{HTML}{E8F3FC}
\definecolor{swRow}{HTML}{F1F8FE}
\definecolor{swSub}{HTML}{DCEBF9}
\definecolor{swRegion}{HTML}{5E8FC2}
\definecolor{swSel}{HTML}{BFDCF7}
\definecolor{swGood}{HTML}{86EFAC}
\definecolor{swMid}{HTML}{FDE047}
\definecolor{swLow}{HTML}{FDBA74}

\pagestyle{fancy}
\fancyhf{}
\fancyhead[L]{\small Kala Care ERP}
\fancyhead[R]{\small PMS -- Performance Management System}
\fancyfoot[C]{\thepage}
\renewcommand{\headrulewidth}{0.4pt}

\fancypagestyle{plain}{%
  \fancyhf{}\fancyfoot[C]{\thepage}\renewcommand{\headrulewidth}{0pt}}

\setlength{\parindent}{0pt}
\setlength{\parskip}{6pt}

\captionsetup{font=small,labelfont=bf}

% ---- helpers ---------------------------------------------------------------
% Requirement list: R1., R2., ... with a hanging indent (same as MOM.pdf)
\newlist{reqlist}{enumerate}{1}
\setlist[reqlist]{label=\textbf{R\arabic*.},leftmargin=3.1em,itemsep=6pt,
                  topsep=6pt,labelsep=0.6em}

\setlist[itemize]{leftmargin=1.6em,itemsep=3pt,topsep=4pt}

% A screenshot that degrades to a framed placeholder when the image is absent.
\newcommand{\screenshot}[2]{%
  \begin{figure}[H]\centering
  \IfFileExists{#1}%
    {\includegraphics[width=\linewidth,keepaspectratio]{#1}}%
    {\setlength{\fboxsep}{0pt}%
     \fbox{\parbox[c][42mm][c]{0.98\linewidth}{\centering\small\itshape
       Screenshot placeholder\\[2pt]\normalfont\ttfamily #1}}}%
  \caption{#2}\end{figure}}

% Section-opening blocks, worded the same way throughout
\newcommand{\covers}[1]{\textit{Covers requirements #1.}\par}
\newcommand{\purpose}[1]{\textbf{Purpose.} #1\par}
\newcommand{\dataused}{\textbf{Data used.}\par}
\newcommand{\devproc}{\textbf{How it works.}\par}
\newcommand{\handle}[1]{\textbf{How to handle.} #1\par}

% A column name of an uploaded file
\newcommand{\col}[1]{\textsc{\small #1}}
% Quoted text
\newcommand{\qt}[1]{\textquotedblleft #1\textquotedblright}
% A colour swatch for the colour-language table
\newcommand{\sw}[1]{\colorbox{#1}{\hspace{1.6em}\rule{0pt}{7pt}}}

\begin{document}

% ============================== TITLE =======================================
\begin{center}
  {\LARGE\bfseries PMS -- Performance Management System}\\[10pt]
  {\large Development Documentation, Kala Care ERP System}\\[14pt]
  {\normalsize Prepared by: Vishal Chavan \& Nikhil Reddy \quad$|$\quad Date: August 2026}
\end{center}

\vspace{10pt}

% ============================== 1. INTRO ====================================
\section{Introduction}

This document describes the \textbf{PMS (Performance Management System)} module of the
Kala Care ERP. PMS is the module that turns the KOEL and ERP data exports into the
management sheets the business actually reviews: branch sales against the annual
operating plan (AOP), service-engineer productivity, SR allocation and closure, the
training master, and the yearly reports -- Service Penetration, AMC \& Bandhan
Projection, Customer Delight Index and Service Load \& Response.

One rule shapes the whole module. Every figure PMS prints is either \textbf{counted}
from an uploaded file or \textbf{entered} by a person in the AOP \& Master page --
never invented. A target nobody has entered prints a dash rather than a zero, and a
figure that no upload can produce honestly (last year's closed actual, FTR / FVR) is
typed by hand and is marked as such in this document. Where an uploaded file's own
column disagrees with something the system could calculate for itself, the file's
column wins, and the reason is written down.

The module is organised into the following pages:
\begin{itemize}
  \item \textbf{AOP \& Master} -- every target and every mapping list the reports read,
  in nine tabs.
  \item \textbf{Sales \& Labour Report} -- Spare and Labour sale against the monthly AOP.
  \item \textbf{Employee Productivity} -- service-engineer-wise SR closure, leads and CDI.
  \item \textbf{SR Allocation Report} -- SRs allocated against SRs closed, engineer wise.
  \item \textbf{Training Report} -- the engineers' skill and training master.
  \item \textbf{Annual Reports} -- one page holding the four yearly sheets.
\end{itemize}

The document is organised as follows: first the roles and all requirements received
from the client; then the files that feed the module and the journey a figure takes
from an Excel file to a report; then the development explained \textbf{page by page and
report by report} -- each section stating what the screen is for, \emph{which data it
uses}, how that data flows into every column, how the page is handled, and the
screenshot of the delivered screen.

% ============================== 2. ROLES ====================================
\section{Roles \& Access}

PMS is not a role, it is a \emph{right}. Any user can be given \qt{PMS Access} from
Profile $\rightarrow$ Edit Employee; the Master Admin always has it. The AOP \& Master
page is controlled more finely still -- \textbf{tab by tab} -- so a person can be given
only the tabs they own, and each of those either as read-only or as editable.

\begin{table}[H]\centering\small
\begin{tabularx}{\linewidth}{@{}p{3.4cm}X@{}}
\toprule
\textbf{Role / Right} & \textbf{Capabilities} \\
\midrule
Master Admin &
Full access to every PMS page and every AOP \& Master tab; uploads and clears data;
maintains the Service Engineer UID master; grants PMS rights to other users. \\
\addlinespace
PMS Access\newline (any role) &
Opens the PMS pages granted to them -- Sales \& Labour, Employee Productivity,
SR Allocation, Training Report and Annual Reports. AOP \& Master opens only if at
least one of its tabs has been granted. \\
\addlinespace
AOP right: \textit{Show only} &
The granted tabs open read-only: the save, sync and reset buttons are hidden and every
cell is locked. The same restriction is applied again by the system itself, tab by tab,
so a read-only right cannot be worked around. \\
\addlinespace
AOP right: \textit{Show and edit} &
The granted tabs can be edited and saved. \\
\addlinespace
Export &
Excel export is available to the Master Admin always, and to other users when the
export permission is enabled on their account. \\
\bottomrule
\end{tabularx}
\caption{Role- and right-based access for the PMS module}
\end{table}

The nine AOP \& Master tabs that rights are granted on are: \textit{Target Master};
\textit{SR Type Master} for Sales and Labour, for MaxTTR, for EFSR and for Service Load;
\textit{Lead Category Master}; \textit{CDI Target Master}; \textit{AMC \& Bandhan AOP};
and \textit{Service Load AOP}.

% ============================== 3. REQUIREMENTS =============================
\section{Client Requirements}

The following requirements were received from the client and finalised for the PMS
module. They are referenced by number in the page-wise sections that follow.

\begin{reqlist}

\item \textbf{Accumulating Spare / Labour upload.} Two Excel files feed the sales side --
the \emph{Part Sale} (spares) file and the \emph{Labour Revenue} file. Rows accumulate
across daily, weekly or monthly extracts, and re-uploading the same file must never
double-count a sale. A changed line updates itself, an identical line is reported as a
duplicate, and the upload always reports how many rows were new, updated, duplicate or
unusable.

\item \textbf{Forgiving column matching.} A file must import even when the header
spelling moves -- \qt{Claim Invoice Date} and \qt{CLAIM  INVOICE DATE.} are one and the
same column. Only three columns are essential: \col{branch id},
\col{claim invoice date} and \col{net taxable amount}. A file missing any of them is
refused, by name, before anything is stored; every other column of the file is kept.

\item \textbf{Cancelled invoices.} An invoice can be cancelled from the preview. The row
stays in the system as a record of what was uploaded but is left out of every report,
and it can be restored.

\item \textbf{AOP Target Master.} Branch-wise \emph{Spare} and \emph{Labour} targets for a
whole financial year (April--March) on one screen, entered and shown in Lakh. The
responsible person shown on the report is the branch's Branch Admin, and the branch list
follows the ERP's own branch master.

\item \textbf{Working-days calendar.} A month's target is spread over its \emph{working}
days, never its calendar days: Sundays are off, and the real holidays are ticked on a
calendar. Maharashtra and Karnataka keep separate calendars, so each region's branches
are measured on their own working days.

\item \textbf{Sales and Labour report for any period.} Month-to-date by default, and any
week, month, quarter or year on the period picker. Every branch shows the period Target,
the Daily Target, the sale on the last day, the Target till date, the Achievement till
date, the invoice count, \% Achieved, Short-Fall and the Balance for the month.

\item \textbf{Independent as-on dates.} The Part Sale and the Labour files usually end on
different dates. Each side must be measured against the targets up to \emph{its own} last
data date, so a labour file that is a week behind never reads as a shortfall.

\item \textbf{Breakdown tables.} The same period broken down Region-wise (with its AOP
target), Segment-wise, Service-Report-Type-wise and Category-wise (spares only, with
quantity and line items).

\item \textbf{Branch-wise report.} One or more branches picked by hand, with Week-wise,
Month-wise, Segment-wise, Service-Report-Type-wise and Category-wise detail worked out
from only those branches' sales.

\item \textbf{Whole-year report.} FY, Quarterly and Month-wise target against achievement
for the full financial year, independent of the period picker, including the run-rate
still needed to close the year.

\item \textbf{SR Type grouping masters.} The raw Service Report Types that arrive in the
uploaded files must be grouped into the reporting heads the business reads. Four separate
grouping lists are needed, because four files carry different vocabularies and two
reports group the \emph{same} file differently. New values appearing in a file must be
picked up automatically for mapping, and a type nobody has mapped must stay visible as
\qt{Unmapped} rather than being quietly merged into something else.

\item \textbf{Service Engineer UID master.} The files name an engineer in two different
ways -- some carry the \col{se name}, others the \col{service engineer uid}. One master
list bridges the two, and matches names ignoring spacing and capitals, so one engineer is
never split into two on the reports. It also lets an engineer whom no file can place be
pinned to a branch by hand.

\item \textbf{Employee Productivity report.} Per service engineer: SRs closed (split by SR
Type), working days, days present, productivity, CDI feedback, leads raised (split by
product category) and the spare and labour conversion amounts -- rolled up by branch,
branch group, region and overall.

\item \textbf{Attendance on task end, not on closure.} Days present must count the days
the engineer actually finished a job in the field, not the day the office closed the SR
in the system -- office closure happens later and in batches.

\item \textbf{Leads attributed through the master only.} A lead counts for an engineer
because its UID is that engineer's UID \emph{in the Service Engineer UID master}. A UID
nobody has mastered is reported as an unlinked lead, never guessed onto a name.

\item \textbf{SR Allocation report.} SRs \emph{allocated} to an engineer and SRs
\emph{closed} by them, engineer wise, with the date columns readable by day, week or
month, plus the branch-wise open SR load.

\item \textbf{Training Report.} The KOEL Training Report export kept as a master --
eight fixed columns and every other column of the file carried along automatically --
searchable two ways: by employee (their skills and full record) and by skill (everyone
who holds it).

\item \textbf{Annual: Service Penetration.} How much of the installed base was actually
touched by a closed service call -- IND and PG population, the unique assets serviced,
and the penetration \%.

\item \textbf{Annual: AMC and Bandhan Projection.} Per branch: last year's actual, this
year's AOP projection and the best month (all entered by hand), the live AMC population
and the month's business done (both counted).

\item \textbf{Annual: AMC monthly sheet.} The same AMC data read by agreement
\emph{category} instead of by branch -- Bandhan with its MH / KAR split, KALA AMC, KOEL
Corporate AMC, expiries, renewals and Live AMC -- one financial year wide, with the
running month opened up into its weeks.

\item \textbf{Annual: Customer Delight Index.} Promotor, Passive and Detractor feedback
per branch, scored as (Promotor $-$ Detractor) $\div$ all feedback $\times$ 100, against
the CDI AOP, with the two previous financial years shown as cumulative columns.

\item \textbf{Annual: Service Load and Response.} SR load by type and by branch,
productivity (calls per person per day), the 4-hour response and the 24 and 48-hour
closure compliance, and the FTR / FVR figures -- as six tabs of one sheet.

\item \textbf{Entered figures where nothing can be counted.} FTR \%, FVR \%, last year's
closed actual and every AOP target are entered by a person in AOP \& Master and printed
exactly as entered. A period nobody has filled in shows a dash, which is not the same
statement as 0\%.

\item \textbf{Excel export everywhere.} Every report exports to a formatted workbook that
mirrors the screen -- the ERP blue band, the header row, the region totals and the
overall row -- and respects the export permission. Dates must export as real dates so
Excel's own date filters keep working.

\item \textbf{Rights per user.} PMS access is granted per user, and the AOP \& Master page
tab by tab, as \emph{none}, \emph{show only} or \emph{show and edit}.

\item \textbf{Reports must open quickly.} The reports run on years of accumulated data,
so a report must summarise its data once and then answer instantly: changing the period,
the day / week / month view or a filter must redraw the screen without waiting for the
data again.

\item \textbf{Branches identified by code, never by name.} Files spell the same branch
differently and the spelling changes between exports, so a branch is always resolved by
its code. A branch code that belongs to no KALA branch -- another dealer's code, which
one of the files carries on about 1\% of its rows -- must collect in a single visible
\qt{Unmapped Branch} row rather than growing a branch that does not exist.

\item \textbf{Consistent India-local dates.} Every date and stamp is India local time, so
a figure reads the same for every user whatever machine they open it on.

\end{reqlist}

% ============================== 4. SOURCE FILES =============================
\section{The Data PMS Reads}

PMS reads thirteen Excel files. Two of them -- the sales files -- are uploaded on the
Sales \& Labour page itself, and the Training file on the Training Report page. Every
other file is uploaded once on the ERP's \textbf{Import} page and is shared with the
rest of the ERP, so PMS never asks for the same file twice.

\begin{table}[H]\centering\small
\begin{tabularx}{\linewidth}{@{}p{4.5cm}p{3.3cm}X@{}}
\toprule
\textbf{File} & \textbf{Uploaded on} & \textbf{Which PMS reports read it} \\
\midrule
Part Sale (spares) & Sales \& Labour & Sales \& Labour report, every view \\
Labour Revenue & Sales \& Labour & Sales \& Labour report, every view \\
Training Report & Training Report & Training Report (By Employee / By Skill) \\
\addlinespace
Response Time \& MaxTTR Details & Import page & Employee Productivity; SR Allocation
(to place an engineer); Service Penetration (assets serviced); Service Load and
Response \\
EFSR Report & Import page & SR Allocation; Employee Productivity (allocated SRs) \\
LMS Data for ERP & Import page & Employee Productivity (leads and conversion amounts) \\
CDI Detail Report & Import page & Employee Productivity (CDI); Annual CDI sheet \\
Asset Detailed Report & Import page & Service Penetration (the installed base) \\
AMC Population Report & Import page & AMC \& Bandhan Projection; the AMC monthly sheet \\
Anubandhan Plus, Anubandhan, BandhanPlus and Regular Bandhan quotes &
Import page & AMC \& Bandhan Projection (the month's business done) \\
Open SR Load Report & Import page & SR Allocation (branch-wise open SR) \\
\bottomrule
\end{tabularx}
\caption{The uploaded files behind the PMS reports}
\end{table}

Two rules apply to every file, whichever page it arrives on. Dates are read as dates
only, and a date written year-first is never mistaken for a day-first one, so
\qt{2026-07-01} can never be read as 7 January. Branch codes are levelled to one form,
so \qt{420435-1}, \qt{420435 1} and \qt{420435\_1} are recognised as the same branch.

% ============================== 5. FLOW =====================================
\section{How the Data Flows}

Every number in PMS makes the same journey. Nothing is calculated in a place the reader
cannot see, and nothing is stored twice.

\vspace{4pt}
\begin{center}\small
\fbox{\parbox[c][10mm][c]{27mm}{\centering Excel file\\ uploaded}}$\;\rightarrow\;$%
\fbox{\parbox[c][10mm][c]{27mm}{\centering Checked\\ and stored}}$\;\rightarrow\;$%
\fbox{\parbox[c][10mm][c]{29mm}{\centering Grouped by the\\ master lists}}$\;\rightarrow\;$%
\fbox{\parbox[c][10mm][c]{29mm}{\centering Measured against\\ the AOP}}$\;\rightarrow\;$%
\fbox{\parbox[c][10mm][c]{22mm}{\centering Report\\ and Excel}}
\end{center}
\vspace{4pt}

\begin{enumerate}[leftmargin=2.2em,itemsep=5pt]
  \item \textbf{The file is uploaded.} Before anything is stored, the file is checked:
  the essential columns must be present, and the user is told by name if one is missing.
  A progress bar runs while a large file is read, so the page is never left blank.

  \item \textbf{Rows are stored once.} Each row is recognised by what makes it unique in
  its own file -- the invoice number and its part or SR line for the sales files; the
  engineer, skill, training date and category for the training file. Uploading the same
  extract again therefore updates what changed and reports the rest as duplicates,
  which is what lets the business upload daily, weekly or monthly without ever
  double-counting.

  \item \textbf{Raw values are grouped by the master lists.} An uploaded row carries raw
  vocabulary -- an SR Type, a \qt{Lead Raised For}, a branch code, an engineer name.
  The AOP \& Master page turns each of those into the wording the report prints: SR Types
  become reporting heads, lead reasons become product categories, engineer names and UIDs
  become one person. Anything not yet mapped collects in a visible \qt{Unmapped} bucket
  instead of disappearing.

  \item \textbf{Counted figures meet entered figures.} What was counted from the files is
  then measured against what the business entered in AOP \& Master -- the monthly Spare
  and Labour targets, the working-days calendar, the CDI and Service Load percentages,
  the AMC projections, and the FTR / FVR figures. Counted and entered figures are never
  mixed up: this document states, for every column, which of the two it is.

  \item \textbf{The report is drawn, and exported.} The screen shows the period asked
  for; changing the period, the day / week / month view or a filter redraws it at once,
  because the report already holds its data. The export writes the same table to Excel.
\end{enumerate}

\textbf{What this means in practice.} A new figure on a report is never a code change
away -- it is usually a mapping the business owns. A new SR Type from KOEL appears in
the master ready to be mapped; a new branch appears from the branch master; a new target
is typed into AOP \& Master. The reports follow.

% ============================== 6. PAGE-WISE ================================
\section{Page-wise Development Process}

This section explains, page by page and report by report, what each screen does, which
client requirements it covers, \textbf{which data it uses}, how that data becomes each
column, how the page is handled, and the screenshot of the delivered screen.

% ---------------------------------------------------------------------------
\subsection{AOP \& Master -- Target Master}
\covers{R4, R5, R27}

\purpose{The annual operating plan itself: every branch's Spare and Labour target for a
whole financial year, and the calendar those targets are spread over.}

\dataused
Nothing is uploaded here -- this tab is where the business \emph{enters} its plan. The
branch list, the branch names and the responsible person come from the ERP's own branch
master, so the AOP grid and the rest of the ERP always name a branch identically.

\devproc
\begin{itemize}
  \item The grid shows a whole financial year at once, April to March, with one row per
  branch and two lines of entry per branch -- Spare and Labour -- and quarter and year
  totals adding themselves up as the user types.
  \item Targets are entered and shown in \textbf{Lakh}. A branch with no target for a
  month leaves an empty cell rather than a wall of zeros.
  \item Working days are set per month \emph{and per region}, because Maharashtra and
  Karnataka do not share holidays. Ticking the real holidays on the calendar makes every
  part-period exact: the working days of any date range are the days in it, less Sundays,
  less the dates ticked for that region.
  \item The order of preference is deliberate. A month with holidays ticked uses the
  calendar, because it knows \emph{which} days were off and not merely how many. A month
  with nothing ticked uses its own typed count; failing that, a standing count for that
  calendar month; failing that, all days except Sundays.
  \item The spelling of a branch entered here is treated as the correct one by every
  other PMS report, so one branch reads the same on every sheet.
\end{itemize}

\handle{Pick the financial year, enter the Spare and Labour targets in Lakh, open a
month's calendar to tick the MH and KA holidays, and Save. Every figure on the Sales \&
Labour report -- Target, Daily Target, Target till date, \% Achieved -- follows from
this tab.}

\screenshot{screens/aop-target-master.png}{AOP \& Master -- Target Master: the
financial-year grid and the working-days calendar}

% ---------------------------------------------------------------------------
\subsection{AOP \& Master -- the four SR Type Masters}
\covers{R11}

\purpose{Turn the raw Service Report Types that arrive in the uploaded files into the
reporting heads each report prints.}

\dataused
The SR Type column of the uploaded files -- and only that column. The tab reads the
values already present in the uploaded data and asks the business what each one means.

There are four of these masters and not one, and the reason matters:

\begin{table}[H]\centering\small
\begin{tabularx}{\linewidth}{@{}p{3.6cm}p{3.5cm}X@{}}
\toprule
\textbf{Tab} & \textbf{Reads which column} & \textbf{Groups them into} \\
\midrule
SR Type Master\newline (Sales and Labour) &
Part Sale \col{claim invoice sr type}; Labour \col{sr type} &
Warranty, Post Warranty, AMC, KOEL AMC, OTC Order \\
\addlinespace
SR Type Master\newline (MaxTTR) &
\col{sr type} of the Response Time \& MaxTTR file &
Warranty, Post Warranty, AMC, KOEL AMC, Courtesy Visit, Others \\
\addlinespace
SR Type Master\newline (EFSR) &
\col{sr type} of the EFSR Report &
the same six; the EFSR file also brings Commercial and RECD Kit \\
\addlinespace
SR Type Master\newline (Service Load) &
the same MaxTTR \col{sr type}, grouped a \emph{second} way &
CSP, Post Warranty, Warranty, KOEL AMC, Dealer AMC, Courtesy Visit, Others \\
\bottomrule
\end{tabularx}
\caption{Why four SR Type masters exist}
\end{table}

\devproc
\begin{itemize}
  \item Each tab lists the SR Types found in \emph{its own} file. \textbf{Sync} re-reads
  the uploaded data and adds any value seen for the first time, so a new type from KOEL
  is never lost -- it simply appears, waiting to be mapped.
  \item \textbf{Reset} restores the mapping the business originally gave. The mapping the
  reports actually read is always the one saved on this screen.
  \item A head can be added or removed from the list beside the mapping, so the reporting
  buckets can change without any development work.
  \item An SR Type with no head is never dropped: it collects in a visible \qt{Unmapped}
  bucket, so the head-wise breakdown always adds up to the total and the gap is obvious.
  \item Sales and Labour, and Service Load, group the same MaxTTR column differently on
  purpose. Employee Productivity folds CSP into Warranty and Dealer AMC into AMC, while
  the Service Load sheet prints CSP and Dealer AMC as rows of their own. Sharing one list
  would force one of the two reports to be wrong.
\end{itemize}

\handle{Open the tab, press Sync after a new upload, give every unmapped type its head,
and Save. The change shows on the next report.}

\screenshot{screens/aop-sr-type-master.png}{AOP \& Master -- SR Type Master: the synced
types and their heads}

% ---------------------------------------------------------------------------
\subsection{AOP \& Master -- Lead Category Master}
\covers{R13, R15}

\purpose{Turn the LMS file's \col{lead raised for} values into the columns of the
Employee Productivity report's \emph{Product Wise Lead Count} group.}

\dataused
The \col{lead raised for} column of the LMS Data for ERP file.

\devproc
\begin{itemize}
  \item The category list is the report's column list, in the order set here -- Allied
  Oil, Allied Product, AMC, Battery, Coolant, DEF, K-Oil, Others, Spares, Whole Goods,
  and anything added later.
  \item Every \col{lead raised for} value in the uploaded data is brought in by Sync and
  given a category by hand -- for example \emph{BD Spares} to \emph{Spares}, and
  \emph{New DG} to \emph{Whole Goods}.
  \item A value with no category collects in an \qt{Unmapped} column, so the product
  split always adds up to the lead count.
\end{itemize}

\handle{Sync after an LMS upload, map the new values, Save.}

\screenshot{screens/aop-lead-category-master.png}{AOP \& Master -- Lead Category Master}

% ---------------------------------------------------------------------------
\subsection{AOP \& Master -- CDI Target Master}
\covers{R21, R23}

\purpose{The AOP column of the Annual Reports' Customer Delight Index sheet -- one
percentage per financial year, per row of that report.}

\dataused
Entered by the business. Nothing is counted here.

\devproc
\begin{itemize}
  \item A CDI score is a percentage, and percentages cannot be added up, so the target is
  set per \emph{row of the report} rather than per branch: a branch row, a region row
  (MH or KA), and the KCGL Overall row each carry their own target.
  \item The printed sheet sets the target at region and overall level only. Branch
  targets are optional and simply stay empty until somebody fills them in -- nothing is
  derived from anything else, and an unset row shows no AOP.
\end{itemize}

\handle{Pick the financial year, type the target \% on the rows the business sets, Save.}

\screenshot{screens/aop-cdi-target.png}{AOP \& Master -- CDI Target Master}

% ---------------------------------------------------------------------------
\subsection{AOP \& Master -- AMC \& Bandhan AOP}
\covers{R19, R20, R23}

\purpose{The figures on the two AMC sheets that the business \emph{states} rather than
counts. The tab carries two tables, because the two sheets have different kinds of row.}

\dataused
Entered by the business. The counted figure for last year is shown alongside for
reference only, and is never printed on the report by itself.

\devproc
\begin{itemize}
  \item \textbf{Per branch} (the first table) -- three columns of the AMC \& Bandhan
  Projection sheet: last year's actual, this year's AOP projection, and the year's best
  month. All three come from here and nowhere else; an empty cell prints a dash rather
  than a figure worked out behind the reader's back.
  \item Last year's actual is entered by hand for a concrete reason. The AMC Population
  Report is a \emph{snapshot}: one row per genset carrying only its latest agreement, so
  a renewal replaces the agreement it renewed and a closed year's count shrinks every
  month. Counted from today's file, FY26 gives 937 against the 1,353 the business
  actually closed it with -- which is why that column is stated, not counted. The screen
  records who entered it and when.
  \item \textbf{Per agreement category} (the second table) -- the AOP of each row of the
  AMC monthly sheet. Those rows are agreement categories, not branches, so their target
  cannot be the sum of branch targets and needs a figure of its own.
  \item Region and company rows need no target of their own here, because these are
  counts and the report simply adds its branches up. (The CDI target is a percentage,
  which is exactly why that one does need a target per row.)
\end{itemize}

\handle{Pick the financial year, fill the branch table and the category table, Save. The
figures appear on the Annual Reports' two AMC tabs.}

\screenshot{screens/aop-amc-bandhan.png}{AOP \& Master -- AMC \& Bandhan AOP: the branch
table and the category table}

% ---------------------------------------------------------------------------
\subsection{AOP \& Master -- Service Load AOP}
\covers{R22, R23}

\purpose{Every target the Service Load and Response sheet is measured against, plus the
two rows nothing uploaded can produce.}

\dataused
Entered by the business.

\devproc
\begin{itemize}
  \item \textbf{SR closure target} -- one number per month, per branch. Everything else
  on that sheet is built from it: a branch row is its twelve months added up, a region
  row is those months across the region's branches, the Total row is the same across
  every branch, and the sheet's monthly \emph{AOP Target} strip is one month across every
  branch. One set of numbers, and every roll-up agrees with it by construction.
  \item \textbf{Percentage and ratio targets} -- Productivity (Calls PP PD), 4 Hrs
  Response, SR Closed within 24 Hrs and SR Closed within 48 Hrs. These cannot be added
  up, so each is set per row: branch, region and overall.
  \item \textbf{FTR and FVR} -- entered per month and per cumulative financial year. The
  year they belong to is fixed, so a figure entered for FY25--26 still means FY25--26
  when the report moves on to the next year. A period nobody has filled in shows a dash,
  which is not the same statement as 0\%.
\end{itemize}

\handle{Enter the monthly SR-closure targets per branch, the percentage AOPs at region
and overall level, and the FTR / FVR figures the business supplies. Save.}

\screenshot{screens/aop-service-load.png}{AOP \& Master -- Service Load AOP}

% ---------------------------------------------------------------------------
\subsection{Profile -- Service Engineer UID Master (Master Admin only)}
\covers{R12, R15, R27}

\purpose{The bridge between the files that \emph{name} a service engineer and the files
that \emph{number} one. Without it, no lead and no allocated SR can be credited to
anybody.}

\dataused
Engineer names and UIDs, read out of the Response Time \& MaxTTR, LMS and EFSR files, and
maintained by hand or by a simple Excel import.

\devproc
\begin{itemize}
  \item The Response Time \& MaxTTR file carries only the engineer's \col{se name}, while
  the LMS and EFSR files identify the same person by \col{service engineer uid}. This
  master holds both, and matches names ignoring spaces and capitals, so a difference in
  spelling between two files never splits one engineer into two on a report.
  \item \textbf{Sync} scans the uploaded files for engineer names and adds the ones the
  master does not have yet, noting which file each was found in. UIDs can also be
  imported in bulk from an Excel sheet.
  \item One person may carry more than one UID -- an engineer who appears under two
  numbers in the source systems is still one person on the reports.
  \item The row can also hold a \textbf{branch}. It is the last resort for an engineer
  whose files give no valid KALA branch code, and is the only way to place such a person
  on a report by hand.
  \item This master stays Master Admin only, whatever other PMS rights a user has.
\end{itemize}

\handle{After a new import, press Sync, fill in the missing UIDs, and pin a branch for
any engineer the reports are showing under \qt{Unmapped Branch}.}

\screenshot{screens/profile-se-uid-master.png}{Profile -- Service Engineer UID Master}

% ---------------------------------------------------------------------------
\subsection{Sales \& Labour Report -- Upload and Preview}
\covers{R1, R2, R3}

\purpose{Get the Part Sale and Labour Revenue files into the system safely, and let the
user see exactly what has been stored.}

\dataused
The two sales files, stored column for column. The layout below is also the order the
preview and the exported workbook use.

\begin{table}[H]\centering\small
\begin{tabularx}{\linewidth}{@{}p{6.5cm}X@{}}
\toprule
\textbf{Part Sale (spares) file} & \textbf{Labour Revenue file} \\
\midrule
ZONE NAME, SD ID, SD NAME, \textbf{BRANCH ID*}, BRANCH NAME, INSTANCE ID, SEGMENT,
PRODUCT SEGMENT, APPLICATION CODE, ENGINE SERIAL NO, CLAIM INVOICE NUMBER,
\textbf{CLAIM INVOICE DATE*}, CLAIM INVOICE SR TYPE, CLAIM INVOICE SR SUB TYPE, CATEGORY,
PART CATEGORY, PART NUMBER, PART DESCTRIPTION, QUANTITY,
\textbf{NET TAXABLE AMOUNT*}
&
ZONE NAME, SD ID, SD NAME, \textbf{BRANCH ID*}, BRANCH NAME, CLAIM INVOICE NUMBER,
\textbf{CLAIM INVOICE DATE*}, PRODUCT SEGMENT, SEGMENT, SERIES, SR NUMBER, SR TYPE,
SR SUBTYPE, \textbf{NET TAXABLE AMOUNT*} \\
\bottomrule
\end{tabularx}
\caption{The two sales files. * marks the three essential columns; every other column is
stored as well, and anything unrecognised is kept rather than dropped}
\end{table}

\devproc
\begin{itemize}
  \item \textbf{Check file format} runs a rehearsal: it reports which columns were found,
  which are missing and a small sample of the data, and stores nothing. The panel also
  lists the full expected column list of each file, with the essential columns marked.
  \item The import shows live progress, so a large file never leaves the page looking
  frozen.
  \item A row is recognised by its invoice number together with the line inside that
  invoice. The Part Sale file has one row per part line; the Labour file has one row per
  labour line -- a single SR is billed as several lines, so the SR number alone is not
  enough to tell two lines apart. This is what lets the same period be uploaded twice
  without the labour figure moving.
  \item Every upload reports how many rows were inserted, updated, duplicate or unusable,
  and the upload itself is kept as a record, so what was loaded and when is always
  answerable.
  \item The preview shows the stored rows in the same column order as the standard Excel
  file, with a search box, and lets an invoice be \textbf{cancelled}: the row stays in
  the system and is left out of every report until it is restored.
\end{itemize}

\handle{Check the format, upload the Part Sale and Labour files -- any period, as often
as needed -- read the preview to confirm, and cancel any invoice the business has
withdrawn.}

\screenshot{screens/sales-upload-preview.png}{Sales \& Labour -- the upload boxes and the
Uploaded File Preview}

% ---------------------------------------------------------------------------
\subsection{Sales \& Labour Report -- All (Spare + Labour)}
\covers{R5, R6, R7, R8, R24}

\purpose{The headline report: every branch's Spare and Labour sale for the chosen
period, measured against the AOP, with four breakdown tables underneath.}

\dataused
\begin{table}[H]\centering\small
\begin{tabularx}{\linewidth}{@{}p{4.4cm}X@{}}
\toprule
\textbf{Where it comes from} & \textbf{What is read} \\
\midrule
Part Sale file & \col{branch id}, \col{claim invoice date}, \col{net taxable amount},
\col{claim invoice number}, \col{zone name}, \col{segment}, \col{product segment},
\col{claim invoice sr type}, \col{category}, \col{quantity} \\
Labour Revenue file & \col{branch id}, \col{claim invoice date},
\col{net taxable amount}, \col{claim invoice number}, \col{zone name}, \col{segment},
\col{product segment}, \col{sr type} \\
AOP \& Master & the monthly Spare and Labour targets, the region of each branch, the
responsible person, and the working-days calendar \\
SR Type Master & the head each SR Type belongs to \\
\bottomrule
\end{tabularx}
\caption{Sales \& Labour -- the data behind the report}
\end{table}

\devproc
\begin{itemize}
  \item The period is month-to-date by default and can be set to any range -- a week, a
  month, a quarter or a year. The Part Sale and Labour sides each keep their \emph{own}
  as-on date, so a file that is a week behind is measured only against the targets up to
  its own last day of data.
  \item Targets are taken from every month the period touches and are shared out over
  \emph{working days} of the branch's own region, so Sundays and the ticked holidays
  carry no target.
  \item Cancelled invoices are excluded everywhere.
\end{itemize}

The branch tables -- one for Spare, one for Labour -- carry these columns, and each
formula is written on the screen itself so a figure never has to be taken on trust:

\begin{table}[H]\centering\small
\begin{tabularx}{\linewidth}{@{}p{4.4cm}X@{}}
\toprule
\textbf{Column} & \textbf{How it is worked out} \\
\midrule
Responsible Person & the branch's Branch Admin, from the Target Master \\
Branch ID / Name / Region & the Target Master, falling back to the uploaded file \\
Target (period) & the full monthly targets of every month the period touches, added up \\
Daily Target & Target $\div$ the working days of the period in that branch's region \\
Achi. on \textit{date} & the sale on the period's last day \\
Target till \textit{date} & each month's target $\times$ its working days inside the
period $\div$ its working days in full \\
Achi. till \textit{date} & the sale over the whole period \\
Invoice Count till & how many different invoices, over the period \\
\% Achieved & Achi. till $\div$ Target till $\times$ 100 \\
Short-Fall & Target till $-$ Achi. till (a negative figure means ahead of plan) \\
Balance / Month & Target $-$ Achi. till \\
\bottomrule
\end{tabularx}
\caption{Sales \& Labour -- the branch table columns}
\end{table}

Four breakdown tables follow, all for the same period:

\begin{itemize}
  \item \textbf{Region-wise Sales} -- Maharashtra first, then Karnataka, in a fixed order,
  each with its own AOP target in front of the sale, since a region's target is simply
  the sum of its branches' targets.
  \item \textbf{Segment-wise Sales} -- read from the row's \col{segment}, falling back to
  \col{product segment}; a row with neither is OTC business.
  \item \textbf{Service Report Type-wise Sales} -- the head each SR Type belongs to.
  Segment and Service Head carry no AOP target, so their \% column is a share of the
  total rather than an achievement against a plan.
  \item \textbf{Category-wise Sales (Spare)} -- spares only, because the Labour file has
  no \col{category} column: sale, quantity, line items and invoice count per category.
\end{itemize}

\handle{Set the period, press Generate, read the two branch tables and then the
breakdowns. Export gives a workbook that mirrors the screen -- Spare and Labour side by
side, region by region.}

\screenshot{screens/sales-report-all.png}{Sales \& Labour -- All (Spare + Labour): the
branch tables and the breakdowns}

% ---------------------------------------------------------------------------
\subsection{Sales \& Labour Report -- Branch-wise Report}
\covers{R9, R24}

\purpose{The same period, but for a hand-picked set of branches, with the detail a
branch review actually needs.}

\dataused
The same two sales files and the same AOP targets as the report above, but read for the
selected branches only.

\devproc
\begin{itemize}
  \item The view opens with every branch selected and narrows as branches are ticked;
  the branch tables above shrink to the selection.
  \item Five detail sections are worked out from \emph{only} those branches' sales:
  \textbf{Week-wise} (the period cut into seven-day stretches, with Spare, Labour and
  invoice count), \textbf{Month-wise}, \textbf{Segment-wise}, \textbf{Service Report
  Type-wise} and \textbf{Category-wise (Spare)}.
  \item Each week also carries its share of the AOP target -- a month's target spread
  over its working days, and that week's days added up -- on the same calendar the main
  report uses, so Sundays and ticked holidays carry no target. When the selection spans
  both regions, a day counts as off only when it is off in every region selected.
\end{itemize}

\handle{Choose Branch-wise, tick the branches, open the sections needed, export.}

\screenshot{screens/sales-report-branchwise.png}{Sales \& Labour -- Branch-wise Report}

% ---------------------------------------------------------------------------
\subsection{Sales \& Labour Report -- FY / Quarterly / Month-wise}
\covers{R10, R24}

\purpose{The whole financial year on one screen: how the year is tracking, quarter by
quarter and month by month, and what run-rate is needed to close it.}

\dataused
The two sales files for the whole financial year, against the full monthly targets from
the Target Master, with the year's working days from the same calendar.

\devproc
\begin{itemize}
  \item These three panels ignore the period picker entirely -- they always cover the
  full financial year asked for, April to March.
  \item The targets used here are the full monthly figures, never shared out over part of
  a month; cancelled rows are excluded, as everywhere else.
  \item The report also knows how far into the year the sales data actually reaches, and
  splits the year's working days into elapsed and remaining, per region. That is what
  makes the run-rate honest -- it rests on the same calendar the targets do.
  \item Quarters are the financial ones: Q1 is April to June, Q4 is January to March.
  \item The export writes three sheets in one workbook -- the FY table, Quarterly and
  Month-wise -- with Spare and Labour across the year.
\end{itemize}

\handle{Pick the financial year, open the FY, Quarterly or Month panel, export.}

\screenshot{screens/sales-report-fy.png}{Sales \& Labour -- FY / Quarterly / Month-wise}

% ---------------------------------------------------------------------------
\subsection{Employee Productivity (Service Engineer Productivity)}
\covers{R12, R13, R14, R15, R24, R27}

\purpose{One row per service engineer answering four questions at once: how much did
they close, how many days did they work, how satisfied were their customers, and how
much new business did they raise.}

\dataused
The \textbf{Response Time \& MaxTTR Details} file is the backbone. Every SR in it that
has a close date and an engineer name counts once for that engineer on that date. The
labour file is deliberately not involved: an SR counts because the engineer closed it,
not because it was billed. Three more files and four master lists hang off that
backbone.

\begin{table}[H]\centering\small
\begin{tabularx}{\linewidth}{@{}p{3.1cm}p{4.3cm}X@{}}
\toprule
\textbf{Group of columns} & \textbf{Source file} & \textbf{What is read} \\
\midrule
Close SR, and its SR Type split &
Response Time \& MaxTTR Details &
\col{se name}, \col{branch id}, \col{branch name}, \col{sr close date}, \col{sr type}
(grouped by the SR Type Master for MaxTTR) \\
\addlinespace
Days present on Task end &
Response Time \& MaxTTR Details &
the different \col{sr task end date} values each engineer has \\
\addlinespace
Allocated SR and its split &
EFSR Report &
\col{service engineer uid}, \col{sd branch code}, \col{task assigned date},
\col{sr type} (grouped by the SR Type Master for EFSR) \\
\addlinespace
CDI: Passive, Detractor, \% &
CDI Detail Report &
\col{x technician name}, \col{activity end date}, \col{cdi category} \\
\addlinespace
Number of Lead, Product Wise Lead Count, Spare / Labour Conv. Amount &
LMS Data for ERP &
\col{lead number}, \col{service engineer uid}, \col{branch id},
\col{lead created date}, \col{lead raised for} (grouped by the Lead Category Master),
\col{part invoice amount}, \col{labour invoice amount} \\
\addlinespace
Working Days &
AOP \& Master &
the month-wise working days, per region, for the days in the chosen period \\
\bottomrule
\end{tabularx}
\caption{Employee Productivity -- every column and where its data comes from}
\end{table}

\devproc
\begin{itemize}
  \item \textbf{Who is one engineer.} A person is recognised by their name ignoring
  spaces and capitals, together with their branch, so \emph{Vijaykumar Jadhav},
  \emph{VijaykumarJadhav} and \emph{VIJAYKUMAR  JADHAV} are always one engineer however
  the files spell them; the best-written spelling seen is the one displayed.
  \item \textbf{Which branch.} The branch is taken from the SR's own branch \emph{code},
  never from the branch name -- files spell the same branch differently and the spelling
  changes between exports, while the code does not. The name is only a label, and the AOP
  master's spelling is used so this report and the Sales report name a branch identically.
  \item \textbf{Leads} are credited strictly through the Service Engineer UID master: the
  UID on a lead is looked up there to find the engineer, and is never matched against
  names. A UID nobody has mastered is reported as an \emph{unlinked lead} so the master
  can be corrected, and each lead counts once however many rows the file carries for it.
  \item \textbf{Allocated SR} counts on the date the SR was \emph{given} to the engineer,
  not on a closing date -- a closing date is empty for exactly the SRs this column exists
  to show, the ones assigned but not yet finished.
  \item \textbf{Days present} counts the different days the engineer \emph{finished a job
  in the field}. The system's closing date is an office event that often lands days later
  and in a batch, so counting it would overstate attendance.
  \item \textbf{Productivity} is Close SR $\div$ \emph{Working Days} -- not divided by
  days present. Days present is attendance; the target is the man-days the AOP makes
  available. Hiding the Working Days column does not change the division.
  \item \textbf{CDI \%} is (Promotor $-$ Detractor) $\div$ all feedback $\times$ 100. The
  screen shows Passive and Detractor; Promotor is still counted and still drives the \%.
  \item \textbf{Working Days and Days present appear on engineer rows only.} Every total
  row -- branch, group, region, Grand Total -- leaves both blank rather than adding
  man-days into a number nobody would read.
  \item \textbf{Roll-ups.} Branches roll into their branch groups, then into Maharashtra
  and Karnataka, then into the Grand Total. A branch code belonging to no KALA branch
  collects in the single \qt{Unmapped Branch} row, so the report still adds up to the
  files.
  \item Changing the period, showing the weekly split or applying a filter redraws the
  screen at once -- the report already holds its data.
\end{itemize}

\handle{Set the period at the top; use the toolbar to show the weekly split, hide
Working Days, or fold the SR Type, CDI and Lead Count bands; then export. Each \% and
count column carries its formula in the heading, so any figure can be traced back.}

\screenshot{screens/employee-productivity.png}{Employee Productivity -- service engineer
rows with the SR Type, CDI and lead bands}

% ---------------------------------------------------------------------------
\subsection{SR Allocation Report}
\covers{R16, R24, R27}

\purpose{What was given to each engineer and what they finished -- and therefore what is
still open.}

\dataused
\begin{table}[H]\centering\small
\begin{tabularx}{\linewidth}{@{}p{3.4cm}p{3.6cm}X@{}}
\toprule
\textbf{Measure} & \textbf{Source file} & \textbf{What is read} \\
\midrule
Allocated Total Task & EFSR Report &
\col{service engineer uid}, \col{task assigned date}, \col{sr type},
\col{sd branch code} \\
Closed Total Task & EFSR Report &
the same rows, counted on \col{task end date} \\
Open SR (branch wise) & Open SR Load Report &
\col{instance id [asset \#]}, which identifies the customer and therefore the branch,
and \col{sr created date} \\
Placing an engineer & Response Time \& MaxTTR &
\col{se name} and \col{branch id}, used only when the EFSR rows cannot place them \\
SE Name / SE UID & Service Engineer UID master & the UID, and the engineer's proper name \\
SR Type rows & SR Type Master (EFSR) & the head each SR Type belongs to \\
\bottomrule
\end{tabularx}
\caption{SR Allocation -- the data behind the report}
\end{table}

\devproc
\begin{itemize}
  \item One SR normally counts \emph{twice}: once on the day it was allocated and once on
  the day the engineer finished it. The two totals differ by the SRs given out inside the
  period but finished outside it, or not finished at all -- and that gap is the open
  backlog.
  \item \textbf{Closed} deliberately means the day the \emph{engineer} finished the job,
  not the back-office closing date, which lands days later and never arrives at all for
  about a third of the tasks the engineer has actually completed.
  \item Engineers are identified through the Service Engineer UID master, so the name and
  UID read exactly the same here as on the Employee Productivity report. A UID the master
  does not know still counts, under the file's own engineer name, so the report always
  adds up to the file; the screen reports how many such rows there are.
  \item \textbf{Which branch an engineer belongs to.} The EFSR file carries the branch
  code of the SR, and on about 1\% of rows that is another dealer's code. Grouping
  straight on it would scatter one engineer across branches that do not exist, so the
  report groups by the \emph{engineer}: every SR of theirs counts in the branch they
  belong to, decided in this order -- the KALA branch they have the most EFSR rows in;
  failing that, the branch the MaxTTR file puts them in; failing that, the branch pinned
  on their master row; failing that, the single \qt{Unmapped Branch} row.
  \item The \textbf{Open SR} grid uses exactly two pieces of the Open SR Load Report and
  counts the whole accumulated table -- it answers \qt{how many SRs are open in this
  branch}, not \qt{how many were in today's upload}.
  \item The date columns can be read by day, week or month, and a control chooses whether
  they carry Allocated, Closed or both; the two totals are always shown.
\end{itemize}

\handle{Set the period, choose the day / week / month view and the date split, open an
engineer's row for the SR-Type-wise breakdown, export.}

\screenshot{screens/sr-allocation.png}{SR Allocation Report -- allocated against closed,
engineer and date wise}

% ---------------------------------------------------------------------------
\subsection{Training Report}
\covers{R17, R24}

\purpose{The service engineers' skill and training master, readable from either end:
\qt{what has this person been trained on} and \qt{who holds this skill}.}

\dataused
The KOEL \textbf{Training Report} export. Eight columns are fixed -- \col{uid no},
\col{employee ticket number}, \col{full name}, \col{occupation}, \col{skill},
\col{branch name}, \col{branch id} and \col{hire date} -- and \emph{every other column
of the file is carried along automatically}, including zone, job title, category,
training dates and the rest, so the export can gain or lose a column without any
development work.

\devproc
\begin{itemize}
  \item Only \col{uid no} and \col{full name} are essential; without them there is no
  person to file a training under.
  \item A training is recognised by the engineer, the skill, the training date and the
  category together. This was proved against the real 407-row export: the engineer and
  skill alone merged 114 rows into each other, and adding the training date still merged
  107, because the same training is recorded under more than one category. All four
  together lose nothing.
  \item It is a \emph{master}, not a period report: rows build up across uploads and a
  re-uploaded row simply updates itself, so the page always shows an engineer's whole
  training history. An engineer with no training yet still gets a row, with the skill
  blank, so the master lists everyone the file carries.
  \item \textbf{By Employee} lists one row per person; clicking opens their skills first,
  then every other detail the file carries about them. \textbf{By Skill} lists one row
  per skill; clicking opens everyone who holds it. Both are instant, and the page decides
  which details describe the \emph{person} and which describe the \emph{training} on its
  own.
  \item The search ignores spacing and capitals, so \qt{nilesh salunke},
  \qt{NILESHSALUNKE} and \qt{Nilesh  Salunke} all find the same person.
\end{itemize}

\handle{Upload the Training Report export whenever KOEL issues a new one, then search by
employee name or by skill and click the row for the full record. Export to Excel as
needed.}

\screenshot{screens/training-report.png}{Training Report -- By Employee and By Skill}

% ---------------------------------------------------------------------------
\subsection{Annual Reports -- the page}
\covers{R18, R19, R20, R21, R22}

\purpose{One page holding the yearly management sheets. A report picker chooses which
sheet is shown; the period picker at the top owns the reporting period for all of them.}

Each sheet reads its own file, so each brings its own range of dates: the period picker
re-sets itself to the financial year of whichever report is open, and offers only the
years that report's data actually covers. The AMC report appears as two tabs of one
report rather than as two entries in the picker, because both tabs read the same file.

% ---------------------------------------------------------------------------
\subsection{Annual Reports -- Service Penetration}
\covers{R18}

\purpose{How much of the installed base was actually touched by a closed service call.}

\dataused
\begin{table}[H]\centering\small
\begin{tabularx}{\linewidth}{@{}p{3.3cm}p{3.4cm}X@{}}
\toprule
\textbf{Block} & \textbf{Source file} & \textbf{What is read} \\
\midrule
Ind / PG Population, Grand Total & Asset Detailed Report &
\col{branch id}, \col{segment} (IND or PG), \col{commissioning date},
\col{asset operational status} \\
UNQ AST SR CLOSED & Response Time \& MaxTTR &
\col{branch id}, \col{segment}, \col{instance id} (or \col{engine serial no}),
\col{sr close date} \\
Pen \% & worked out & serviced $\div$ population $\times$ 100 \\
\bottomrule
\end{tabularx}
\caption{Service Penetration -- the data behind the report}
\end{table}

\devproc
\begin{itemize}
  \item The population is \emph{cumulative}: each asset counts once, in its own branch,
  on the day it was commissioned, and the sheet reads the installed base \emph{as on} the
  end of the period -- not what happened to be commissioned inside it.
  \item Assets the ERP has retired are out of the base; an asset with no commissioning
  date cannot be placed on a calendar at all and is reported rather than quietly dropped.
  \item \textbf{Unique means unique}: an asset with five closed SRs counts once. The
  report is built so that this stays exact for any period asked for.
  \item An asset is counted under the branch and segment its \emph{own service record}
  names, so nothing is lost when the asset master has not caught up with a new machine
  yet.
  \item Only the Total Pen \% cell is shaded, against a threshold that is deliberately
  not printed anywhere on the sheet -- it decides a colour, it is not a figure the sheet
  reports.
\end{itemize}

\handle{Pick the period; read population, assets serviced and Pen \% per branch, with the
Maharashtra, Karnataka and overall totals. The export mirrors the screen.}

\screenshot{screens/annual-service-penetration.png}{Annual Reports -- Service Penetration}

% ---------------------------------------------------------------------------
\subsection{Annual Reports -- AMC \& Bandhan Projection (tab 1)}
\covers{R19, R23}

\purpose{Per branch: the year's AMC position against the plan, and the business actually
done in the month.}

\dataused
\begin{table}[H]\centering\small
\begin{tabularx}{\linewidth}{@{}p{4.6cm}p{3.4cm}X@{}}
\toprule
\textbf{Column} & \textbf{Where it comes from} & \textbf{How it is produced} \\
\midrule
F\textit{nn} ACT D/BAMC & AOP \& Master & entered -- last year's actual \\
F\textit{nn} PROJ AOP D/BAMC & AOP \& Master & entered -- this year's projection \\
BEST ACT AOP D/BAMC (M) & AOP \& Master & entered -- the year's best month \\
F\textit{nn} YTD ACT AMC NOS & AMC Population Report &
counted: agreements that are still active \emph{and} have not yet expired. This is a
population as on today, so neither the year nor the month picked moves it \\
\textit{Mon}-\textit{yy} Act & the four Bandhan quote files &
counted: the payments received in that month, branch wise \\
\bottomrule
\end{tabularx}
\caption{AMC \& Bandhan Projection -- where each column comes from}
\end{table}

\devproc
\begin{itemize}
  \item The sheet asks two different questions and no single file answers both, so it
  reads two sources: the agreement population, and the money actually received.
  \item \textbf{The month's business is payments, not cover starting.} Counted on the
  payment date, the figures agree with the business's own sheet -- Mar-26 145, Apr-26
  105, Jun-26 104 against 144, 108 and 106 -- whereas counting when the agreements start
  is 15 to 30 out.
  \item \textbf{The first three columns belong to the AOP master and to nothing else.}
  Last year's actual could not be counted honestly anyway: the AMC file keeps only the
  latest agreement for each machine, so a renewal replaces the agreement it renewed and a
  closed year's count shrinks every month.
  \item \textbf{Best month} follows the monthly payments, and a region or company row
  takes the best month of \emph{its own} monthly totals -- never the sum of its branches'
  best months, which fall in different months and would invent a month nobody had.
  \item Every other roll-up is a plain sum of the branches. The entered figures stay
  empty until at least one branch has one, so an unfilled region reads as empty rather
  than as a target of zero.
  \item The month picker moves the last column only. Everything else belongs to the year,
  not to the month being read.
\end{itemize}

\handle{Open the AMC report, stay on the Projection tab, choose how far into the year to
read, export.}

\screenshot{screens/annual-amc-projection.png}{Annual Reports -- AMC \& Bandhan Projection}

% ---------------------------------------------------------------------------
\subsection{Annual Reports -- AMC (tab 2, the monthly sheet)}
\covers{R20, R23}

\purpose{The same AMC data read by agreement \emph{category} instead of by branch, one
financial year wide.}

\dataused
The AMC Population Report -- \col{branch id} (for the MH / KAR split),
\col{agreement type}, \col{agreement start date}, \col{agreement end date},
\col{agreement status} and \col{last agreement number} -- plus the category AOP entered
in AOP \& Master.

\begin{table}[H]\centering\small
\begin{tabularx}{\linewidth}{@{}p{5.4cm}X@{}}
\toprule
\textbf{Row} & \textbf{What it counts} \\
\midrule
KOEL Bandhan Total & the MH and KAR rows added up \\
KOEL Bandhan MH / KAR & the KOEL Bandhan family, split on the branch's region \\
KALA AMC & the dealership's own AMC \\
KOEL Corporate AMC & KOEL's own agreement, printed as a running total \\
AMC Expired During the Month & agreements whose cover ended in that month \\
AMC Renewed During the Month & agreements that replaced an earlier one \\
Live AMC (KOEL+KALA) OVERALL & KOEL Bandhan plus KALA AMC \\
\bottomrule
\end{tabularx}
\caption{The AMC monthly sheet's rows}
\end{table}

\devproc
\begin{itemize}
  \item The columns are: the AOP for the year, the year so far, one column per month
  reached, and the running month \emph{also} broken into its Monday-to-Sunday weeks --
  beside its month column, not instead of it, which is how the business prints the sheet.
  A finished financial year reads as twelve month columns and no weeks.
  \item \textbf{Closed years are deliberately not shown.} The AMC file keeps only the
  latest agreement for each machine, so a closed year's count leaks away as renewals
  replace it: counted from today's file the Bandhan rows give 235 for FY24--25 and 858
  for FY25--26 against the 837 and 1,302 the business closed those years with. A column
  that wrong is worse than no column, so the sheet prints only the running year, which is
  counted live and is sound.
  \item An agreement type nobody has accounted for is reported rather than dropped, so a
  new product from KOEL announces itself instead of vanishing.
  \item The second table, \emph{KCGL Total AMC}, plans against the sum of every branch's
  projection from the AOP master, spread evenly across the twelve months -- so the yearly
  plan lives in exactly one place, per branch where it is actually owned, and the company
  row follows from it instead of being typed a second time and left to drift.
\end{itemize}

\handle{Open the AMC report and switch to the AMC tab. Its AOP comes from AOP \& Master
$\rightarrow$ AMC \& Bandhan AOP, the category table.}

\screenshot{screens/annual-amc-monthly.png}{Annual Reports -- the AMC monthly sheet}

% ---------------------------------------------------------------------------
\subsection{Annual Reports -- Customer Delight Index (CDI)}
\covers{R21, R23}

\purpose{The customer feedback score per branch, against its AOP, for the financial
year.}

\dataused
The CDI Detail Report -- \col{branch name}, \col{activity end date} and
\col{cdi category} -- and the target entered in the CDI Target Master.

\devproc
\begin{itemize}
  \item Every feedback row falls into one of three buckets, read from the \emph{first
  word} of its category so the score range in brackets can change without breaking
  anything: Promotor, Detractor, and everything else Passive.
  \item The score of any set of rows is (Promotor $-$ Detractor) $\div$ (Promotor +
  Passive + Detractor) $\times$ 100. It is therefore \textbf{not} an average of the rows
  beneath it: every total row is scored from its own feedback, which is why a region row
  can sit outside the range of its own branches.
  \item The CDI file is the only one that carries no branch code -- it names its branch
  (\qt{KALA Care Global LLP - Ahmednagar}) and nothing else. The list matching those names
  to branches is therefore built \emph{from the other files}, which carry both the name
  and the code, so when the ERP renames a branch the match follows on the next upload,
  with no development work. The AOP master's own spelling is registered against the same
  branch as well, so either wording is understood.
  \item Feedback for a branch nobody knows collects in the \qt{Unmapped Branch} row, and
  feedback with no activity date cannot be placed on a calendar at all -- both are
  reported on the screen.
  \item The columns follow the financial year the period falls in: the two previous
  financial years as cumulative columns, the AOP, the year so far, one column per whole
  month, and the running month opened up into its Monday-to-Sunday weeks.
\end{itemize}

\handle{Pick the financial year, read branch, region and overall scores against the AOP,
export.}

\screenshot{screens/annual-cdi.png}{Annual Reports -- Customer Delight Index}

% ---------------------------------------------------------------------------
\subsection{Annual Reports -- Service Load and Response}
\covers{R22, R23, R24}

\purpose{The service department's own scorecard: how much load was closed, how quickly it
was responded to, and how productive the engineers were.}

\dataused
\textbf{One file only} -- Response Time \& MaxTTR Details -- read on \col{sr close date}
per \col{branch id}, using \col{sr type}, \col{response time} and \col{se name}. Three
sets of entered figures sit on top of it: the Service Load SR Type grouping, the SR
closure and percentage AOPs, and the FTR / FVR figures.

The printed sheet is four pages; the screen is \textbf{six tabs}, because its response
page carries two independent measures and its closing blocks summarise the whole sheet:

\begin{table}[H]\centering\small
\begin{tabularx}{\linewidth}{@{}p{3.9cm}X@{}}
\toprule
\textbf{Tab} & \textbf{What it shows} \\
\midrule
Service Load & the SR-type-wise breakdown and its Total, then the same load by branch
under MH / KA, closing on the Grand Total \\
Productivity & Calls PP PD per region, branch and overall, plus the
\emph{Service Request Closure (Nos.)} month strip \\
4 Hrs Response & the 4-hour response \% per region and branch \\
SR Closed 24 Hrs & the 24-hour closure \% per region and branch \\
SR Closed 48 Hrs & the 48-hour closure \% per region and branch \\
Overall and FTR / FVR & the MAX TTR Overall block, and the FTR / FVR figures entered in
AOP \& Master \\
\bottomrule
\end{tabularx}
\caption{Service Load and Response -- the six tabs}
\end{table}

\devproc
\begin{itemize}
  \item \emph{Total Service Load Available} is every SR of the period whatever its type;
  the breakdown rows are that same population split by SR Type. A type nobody has grouped
  still counts in the total but gets no breakdown row, so the difference is visible on
  the sheet instead of being swallowed.
  \item \textbf{All three compliance percentages read one figure} -- the file's own
  response time -- over every SR of the period: within 4, within 24 and within 48 hours.
  \item \textbf{The SR open date is never used}, and this is the most important decision
  on the sheet. In the live file, 8,721 of 45,279 rows carry an open date that falls
  \emph{after} their close date, always at midnight: those are the \emph{scheduled} dates
  of planned work that was completed early. So response and turnaround are read from the
  file's own columns and are never re-calculated from open to close, which would be
  meaningless on 42\% of the file.
  \item \textbf{A negative response time counts as zero hours.} The job was closed before
  it was even due, so it is inside every window -- which is exactly how the source file
  scores it. Counting them as zero rather than discarding them keeps the percentages
  describing the same population as the SR total; a cell says on hover how many of its
  SRs those were.
  \item \textbf{Every percentage row is scored from its own SRs}, never averaged from the
  rows beneath it -- a region row divides the region's compliant SRs by the region's SRs.
  Productivity works the same way: a region's closures over the region's man-days.
  \item \textbf{Productivity} is closures $\div$ (working days $\times$ the engineers who
  actually closed something in that window).
  \item \textbf{FTR and FVR are not counted.} First Time Right and First Visit Report are
  judgements KOEL makes about a job, and the two files that once carried them were
  withdrawn -- so the sheet prints what a person entered in AOP \& Master, and a period
  nobody has filled in shows a dash.
\end{itemize}

\handle{Pick the financial year, switch tabs for load, productivity and each response
window, and read the AOP column beside every measure. The export writes the tab you are
on.}

\screenshot{screens/annual-service-load.png}{Annual Reports -- Service Load and Response}

% ============================== 7. CROSS-CUTTING ============================
\section{Cross-cutting Behaviour}

\subsection{Access rights}
\covers{R25}
PMS access is granted per user, and the AOP \& Master page tab by tab, as \emph{none},
\emph{show only} or \emph{show and edit}. On a show-only tab the save, sync and reset
buttons are hidden and every cell is locked -- and the same restriction is applied again
by the system itself, tab by tab, so a read-only right cannot be worked around. The
Service Engineer UID master stays Master Admin only, whatever other PMS rights a user
holds.

\subsection{Speed, and how fresh the figures are}
\covers{R26}
The reports run on years of accumulated data, so each one summarises its data once and
then answers instantly: changing the period, the day / week / month view or a filter
redraws the screen without fetching anything again. A summary is re-used until new data
is uploaded, and a new upload makes the next report rebuild itself -- so a figure on
screen is always as current as the last upload, and never older.

\subsection{Nothing is dropped silently}
Every report shows its own gaps rather than hiding them. SR Types and lead reasons nobody
has mapped get a visible \qt{Unmapped} bucket; branch codes belonging to no KALA branch
collect in one \qt{Unmapped Branch} row; a lead whose engineer is not in the master is
counted as unlinked and reported; rows with no usable date are reported as such. A figure
the system cannot produce honestly is entered by a person and is marked in this document
as entered.

\subsection{Excel export}
\covers{R24}
Every report exports to a workbook that mirrors the screen: the ERP blue band carrying
the report name and its period, the header row, the branch rows, the Maharashtra and
Karnataka totals, and the overall row. The workbook always uses the light colours, since
it is read and printed on its own. Dates are written as real dates, so Excel's own date
filters work on the exported sheet. Export follows the export permission on the user's
account.

\subsection{The colour language of the reports}
Every PMS grid speaks the same colours, on screen and in the exported workbook, so a
reader who learns one report can read them all.

\begin{table}[H]\centering\small
\begin{tabularx}{\linewidth}{@{}p{1.6cm}p{4.2cm}X@{}}
\toprule
\textbf{Colour} & \textbf{Where it appears} & \textbf{What it means} \\
\midrule
\sw{swBand} & the top band and the Grand Total row & the report's own identity and its
final total \\
\sw{swHead} & every column heading & a heading, never a figure \\
\sw{swRow} & branch and engineer rows & an ordinary data row \\
\sw{swSub} & Sub Total and group rows & a group heading the rows below it \\
\sw{swRegion} & Maharashtra / Karnataka rows & a region total \\
\sw{swSel} & the row that has been opened or picked & where the reader is \\
\sw{swGood} & a \% cell & at or above target \\
\sw{swMid} & a \% cell & within a fifth of the target \\
\sw{swLow} & a \% cell & further below the target \\
\bottomrule
\end{tabularx}
\caption{The shared PMS colour language}
\end{table}

A dash is never coloured: an absent figure stays neutral, because \qt{not entered} is not
a judgement about performance. The same colours are used in the dark theme, adjusted for
a dark background.

\subsection{India-local dates}
\covers{R28}
Every date and every stamp in the module is India local time, so a figure reads the same
for every user, on any machine, wherever it is opened.

% ============================== APPENDIX ====================================
\appendix
\section{Glossary}

\begin{table}[H]\centering\small
\begin{tabularx}{\linewidth}{@{}p{3.5cm}X@{}}
\toprule
\textbf{Term} & \textbf{Meaning} \\
\midrule
AOP & Annual Operating Plan -- the targets the business sets for the financial year \\
FY & Financial year, April to March. FY26--27 starts in April 2026 \\
SR & Service Request -- one service job \\
SR Type & What kind of job it was (Warranty, Post Warranty, AMC, CSP and so on) \\
Head & The reporting bucket several SR Types are grouped into \\
Spare / Labour & The two halves of service revenue: parts sold, and labour billed \\
OTC & Over-the-counter business -- a sale with no segment of its own \\
Segment (IND / PG) & The two segments of the installed base the reports measure \\
MH / KA (KAR) & The Maharashtra and Karnataka regions \\
CDI & Customer Delight Index -- the customer feedback score \\
Promotor / Passive / Detractor & The three feedback buckets a CDI response falls into \\
AMC & Annual Maintenance Contract \\
D/BAMC & Dealer AMC plus Bandhan AMC, counted together \\
Live AMC & Agreements that are active and have not yet expired \\
MaxTTR & Maximum turnaround time -- the response and closure timings file \\
EFSR & The report of SRs allocated to and completed by each engineer \\
LMS & The lead management data -- leads raised and what they converted into \\
FTR / FVR & First Time Right / First Visit Report -- quality judgements made by KOEL \\
Calls PP PD & Calls per person per day -- the productivity ratio \\
Working days & The days a target is spread over: all days less Sundays and holidays \\
Unmapped & A value nobody has classified yet -- shown, never hidden \\
\bottomrule
\end{tabularx}
\caption{Terms used in this document}
\end{table}

\vspace{12pt}
\begin{center}\small\itshape End of document\end{center}

\end{document}
